\chapter{Sélection et Visualisation}

Dans ce chapitre, nous expliquons quelles caractéristiques (features) 
sont conservées, et quelles méthodes ont été utiliser pour les choisir. Nous abordons 
également l'étape de Visualisation de nos données.

\section{Sélection}
\subsection{RFECV}

\subsection{ACP}

\subsection{SelectKBest}

% \section{f\_classif}

% \section{mutual\_info\_classif}


\section{Visualisation}
Suite à cette étape de sélections, nous avons choisis de conserver les déscripteurs communs 
aux trois méthode RFECV, ANOVA et mutual information, ansin 20 déscripteurs sont 
conservées pour les modèles de classification: 

\noindent ['FT\_Adef', 'FT\_Aforme', 'FT\_Elo\_A', 'FT\_Elo\_H', 'FT\_Elo\_dif', 'FT\_Elo\_ratio', 'FT\_Hforme', 'FT\_forme\_diff',
'FT\_forme\_ratio', 'FT\_prec\_diff', 'FT\_prec\_ratio', 'HT\_Elo\_A', 'HT\_Elo\_H', 'HT\_Elo\_dif', 'HT\_Elo\_ratio',
'HT\_Hforme', 'HT\_def\_ratio', 'HT\_forme\_diff', 'HT\_forme\_ratio', 'Rank\_diff'].
\newline 
Une liste de 6 déscripteurs sera sélectionnés pour les réseaux de neurones : 

\noindent ['FT\_Aprec\_weight', 'FT\_avgY\_diff', 'HT\_Elo\_A', 'HT\_Elo\_dif', 'HT\_Hdef', 'WinStreak\_diff'].

A présent, nous allons aborder la Visualisation des données par rapport aux deux listes de 
déscripteurs. Nous avons utilisé Umap, ACP biplot ainsi que Boxplot afin de projeter 
differement.

\subsection{Umap}

Nous avons utilisé la méthode Umap afin de réduire la dimensionnalité de nos données. 
Grâce à cette méthode, nous avons projeté nos données en 2 et 3 dimensions.

\begin{figure}[H]
    \centering
    % Première mini-page (gauche)
    \begin{minipage}{0.48\textwidth}
        \centering
        \includegraphics[width=0.9\textwidth]{../../../plot/umap_2d_ft_co.png}
        \caption{UMAP 2D : 20 features communes.}
        \label{fig:umap_20}
    \end{minipage}
    \hfill 
    % Seconde mini-page (droite)
    \begin{minipage}{0.48\textwidth}
        \centering
        \includegraphics[width=0.9\textwidth]{../../../plot/umap_2d_ft_ecd.png}
        \caption{UMAP 2D : 6 features (modèle ECD).}
        \label{fig:umap_6}
    \end{minipage}
\end{figure}

Sur la figure \ref{fig:umap_20} nous avons projeté nos données en utilisant les 
déscripteurs sélectionnés par RFECV, ANOVA, Mutual information. Au total 20 features 
sont retenues communes à ces trois méthodes. Sur la figure \ref{fig:umap_6} nous avons 
projeté en utilisant les déscripteurs retenues grâce au MLP utilisant la méthode ECD 
(basé sur \cite{leray}).

\begin{figure}[H]
    \centering
    % Première mini-page (gauche)
    \begin{minipage}{0.48\textwidth}
        \centering
        \includegraphics[width=\textwidth]{../../../plot/umap_3d_ft_co.png}
        \caption{UMAP 2D : 20 features communes.}
        \label{fig:umap_3d_20}
    \end{minipage}
    \hfill 
    % Seconde mini-page (droite)
    \begin{minipage}{0.48\textwidth}
        \centering
        \includegraphics[width=\textwidth]{../../../plot/umap_3d_ft_ecd.png}
        \caption{UMAP 2D : 6 features (modèle ECD).}
        \label{fig:umap_3d_6}
    \end{minipage}
\end{figure}
Ici, nous avons projeté en 3 dimensions afin de pouvoir mieux explorer la répartition et 
la séparabilité de nos données.


\subsection{ACP}
\begin{figure}[H]
    \centering
    % Première mini-page (gauche)
    \begin{minipage}{0.48\textwidth}
        \centering
        \includegraphics[width=\textwidth]{../../../plot/acp_2d_ft_co.png}
        \caption{acp 2D : 20 features communes.}
        \label{fig:acpd_2d_20}
    \end{minipage}
    \hfill 
    % Seconde mini-page (droite)
    \begin{minipage}{0.48\textwidth}
        \centering
        \includegraphics[width=\textwidth]{../../../plot/acp_2d_ft_ecd.png}
        \caption{acp 2D : 6 features (modèle ECD).}
        \label{fig:acp_3d_6}
    \end{minipage}
\end{figure}



\subsection{Boxplot}

\begin{figure}[H]    
    \centering
    \includegraphics[width=\textwidth]{../../../plot/boxplot_Hvs_ft_co.png}
    \caption{Boxplot : 20 features communes.}
    \label{fig:Boxplot_20}
\end{figure}

Nous avons également comparé la séparabilité des classes en utilisant la cible 'Hvs' 
( Homve versus Rest ). 

\begin{figure}[H] 
    \centering
    \includegraphics[width=\textwidth]{../../../plot/boxplot_Hvs_ft_ecd.png}
    \caption{Boxplot : 6 features (modèle ECD).}
    \label{fig:Boxplot_6}
\end{figure}


Les figures \ref{fig:Boxplot_20} et \ref{fig:Boxplot_6} montre la différence entre 
les features sélectionnés.